The study of cloud coverage and its physical properties are the first step in understanding cloud radiative effects.  In order to retrieve these variables, CLOUDNET post-processed data of FMCW 94 GHz cloud Doppler radar,  HATPRO radiometer and CHM15k ceilometer were used. Additionally, complementary analysis of thermodynamic conditions was performed with the HATPRO radiometer vertical profiles. Temperature and relative humidity profiles tracked cold (fall-winter) and warm regimes (spring-summer).  
The assortment algorithm was applied using the CLOUDNET target classification product to classify different single layer cloud types. Inter-annual observation on cloud occurrence revealed a seasonal trend, where mixed phase clouds are the prevailing species, followed by ice clouds, especially during spring and winter. In these seasons, precipitating mixed phase clouds can reach 30\%XXX and 25\%XXX of frequency, respectively, and ice clouds 15\%XXX and 12\%XXX. Other cloud types have no strong seasonal pattern, and precipitating liquid clouds are absent throughout the year. In addition, the total single layer cloud frequency in summer is less than 8\%.
Mixed-phase clouds exhibit the thickest layers, followed by ice clouds and then liquid clouds. The median cloud thickness for these clouds is 2.0XXX, 1.2XXX, and 0.2XXX, respectively. However, due to their considerable variability, no distinct seasonal pattern was discernible for this property. Consequently, thicker layers of mixed-phase clouds contain more ice content than ice clouds layers. Although typical values for both are 0-20\,g~m$^{-2}$, mixed phase clouds have systematically higher IWP, with an expressive increase in fall, reaching 20-50\,g~m$^{-2}$. Conversely, liquid clouds contain slightly more liquid content than mixed-phase clouds, despite being thin. This indicates that liquid layers consistently maintain thinness, even within mixed-phase clouds. In fact, they have almost the same amount of water, within the interval of 0-20\,g~m$^{-2}$, but liquid clouds have higher probability to reach higher values, specially in spring and fall (see figure \ref{fig:LWP_hist}).

LWC and droplets effective radius does not show a clear seasonal trend for mixed and liquid phase clouds. Effective radius has almost the same distribution along year, with a median of 5\,$\mu$m for liquid clouds, and is found mainly below 3\,km. Differently mixed phase clouds have a radius that can vary a lot from the surface to 10 km, but in general, most frequent values between 5 and 12\,$\mu$m are found between 1-2 \,km and 5-9\,km. 
Concerning LWC, despite the fact that it was easy to identify higher values in cold regimes and smaller ones during summer, the variability of data goes through several orders of magnitude. It was observed for liquid and mixed phase clouds. 

With respect to ice and mixed phase clouds, the same difficulty to reach a seasonal pattern for IWC was identified. Month profiles did not show a significant difference, except that in summer, values are shifted upwards in the atmosphere. For ice clouds, effective radius can reach 65\,$\mu$m in summer, but normally they vary between 20-30\,$\mu$m around 8-12\,km.  For mixed phase clouds, ice effective radius fluctuates around 50\,$\mu$m between 1-7\,km.

Concerning cloud base height, a significantly seasonal difference on vertical distribution was observed in summer. Normally, the median value of ice cloud base height is 8\,kmXXX in fall and 7\,kmXXX for the rest of the year. For mixed phase, median values are 4\,km and 5\,km in cold and warm regimes, respectively. However, an unnusual mode around 5\,km appeared in summer, especially for mixed phase clouds, which also present a sharper distribution compared to other seasons. Nonetheless, in this season, liquid clouds presented a spreader distribution, with a median of 3\,kmXXX, against a median of 1\,kmXXX in other seasons. Cloud top height followed a simmilar behavior. In colder regimes (winter-spring), cloud top is 1.5\,kmXXX, 9\,kmXXX, 6\,kmXXX, for 
liquid, ice and mixed clouds, respectively. During warm periods (summer-fall), medians changed to 2.5\,kmXXX, 8.5\,kmXXX, 7.5\,kmXXX. However, only in summer, the meadian cloud base height for liquid clouds were 3\,kmXXX. Finally, a comprehensive analysis of cloud occurrence and physical properties was conducted at the ACTRIS-CCRESS Granada station. Notably, this is the first study on cloud properties in the southeast of Spain, an area recognized as a hotspot for climate variability. Consequently, it has the potential to enhance our understanding of global cloud properties, as the methods employed here can be applied to other stations as well.
